\section{Compact Appendix E: no second same-domain equivalence and no hidden inferential locus}
\label{app:no-second-equivalence}

This appendix records the final closure step needed by the body: once the kernel, quotient, and coherent-trajectory layers are fixed, there is no second closure-relevant same-domain equivalence relation left to do governance work.

\begin{theorem}[Standing conservation as an identity-level invariant]
Standing conservation is not merely a policy imposed on top of the regime. It is an identity-level invariant of admissible classification on the fixed domain.
\end{theorem}

\begin{proof}
If standing could vary independently of admissible identity, then two states or lineages identical for same-domain admissibility purposes would still differ in standing-bearing force. That would introduce a second closure-relevant classification layer beyond the standing quotient and coherent-trajectory structure already fixed by the kernel. The fixed regime excludes that possibility.
\end{proof}

\begin{theorem}[No second closure-relevant same-domain equivalence]
There is no second same-domain equivalence relation, finer or coarser than the standing quotient together with coherent legal trajectory structure, that carries additional closure-relevant governance force while preserving the fixed kernel.
\end{theorem}

\begin{proof}
A finer relation would distinguish states or lineages that the quotient and coherent-trajectory layers already identify as governance-equivalent, thereby reintroducing hidden standing-bearing structure beneath tensor content or beneath pointwise standing-preservation. A coarser relation would identify states or lineages whose standing or coherent-trajectory behavior still differs and would therefore erase genuine inferential-life distinctions. Neither move preserves the kernel. So no second closure-relevant equivalence survives.
\end{proof}

\begin{corollary}[No hidden inferential locus]
Any same-domain distinction that changes inferential life must surface in one of two and only two places: standing classification of evaluated acts or coherent legal trajectory structure. If it surfaces in neither, it is closure-inert inside the regime.
\end{corollary}

\begin{proof}
The quotient layer fixes all same-domain evaluated standing distinctions, and the coherent-reasoning layer fixes all closure-sensitive same-domain path distinctions. The previous theorem shows that no further closure-relevant equivalence remains. So there is no third locus for same-domain governance-relevant difference to hide.
\end{proof}

\begin{remark}[Formalization and audit placement]
A companion Lean4 development is evidential and audit-supportive only. The primary manuscript relies on the prose theorem chain. Detailed theorem-ledger tables, dependency graphs, and hostile-proof provenance are intentionally offloaded to the audit companion so that the visible manuscript remains application-driven and kernel-relative.
\end{remark}
